\section{Solutions for Chapter~\ref{ch:characters}}

\paragraph{Exercise~\ref{ex:char-oscillator}.}
\phantomsection
\label{sol:ex:char-oscillator}
The five partitions give, respectively,
\[
 8,\quad8^2,\quad\binom92,\quad
 8\binom92,\quad\binom{11}{4}.
\]
Their sum is $8+64+36+288+330=726$.
The $E_8$ theta coefficients through weight four are
$1,240,2160,6720,17520$. Hence the VOA coefficient is
\[
 17520+6720(8)+2160(44)+240(192)+726=213126.
\]
This separates momentum degeneracy from oscillator degeneracy.



\paragraph{Exercise~\ref{ex:char-primary}.}
\phantomsection
\label{sol:ex:char-primary}
Through grade three, multiplication by
$P(q)=1-q-q^2+\cdots$ gives
\[
\begin{array}{c|rrrr}
 &v(0)&v(1)&v(2)&v(3)\\\hline
 b&80&1170&9300&53162\\
 d_4&5&224&2504&16975
\end{array}
\]
For example, $v_b(3)=65042-10630-1250=53162$.
Convolution with $1/P(q)=1+q+2q^2+3q^3+\cdots$
returns
$a(1)=v(1)+v(0)$,
$a(2)=v(2)+v(1)+2v(0)$, and
$a(3)=v(3)+v(2)+2v(1)+3v(0)$.
These recover all the displayed character entries.



\paragraph{Exercise~\ref{ex:char-pole}.}
\phantomsection
\label{sol:ex:char-pole}
A twist determines $h$ modulo one, so the three stated
weights have the same twist. The possible $2/13$ top would
produce a pole with coefficient $\gamma$. Existence of the
specified fundamental modular solutions, together with
$\dim U_1=6+13\gamma$ and the independent six-current theorem,
forces $\gamma=0$. The uniquely selected fundamental solution
then has its first nonzero $d_3$ coefficient at $15/13$.
Excluding $2/13$ by itself would not exclude $28/13$;
that final selection uses the known fundamental solution.



\paragraph{Exercise~\ref{ex:char-equal}.}
\phantomsection
\label{sol:ex:char-equal}
The Grothendieck group is free on the isomorphism classes
of simples. Taking a character is a linear map from that group
to a space of functions, and the map need not be injective.
Defect insertions, including tube projectors, can distinguish
equal unflavoured characters. Flavoured traces or matrix
elements of specified current zero modes can also distinguish
modules when the associated actions differ.



\paragraph{Exercise~\ref{ex:char-uniqueness-scope}.}
\phantomsection
\label{sol:ex:char-uniqueness-scope}
If $v$ is a modular-invariant constant vector, adding $tv$
does not change any negative-power coefficient. The compact
modular-curve argument then proves uniqueness only modulo
such invariant constants. One must prescribe the projection
of the constant coefficient onto the invariant subspace,
or an equivalent value normalization, to remove the ambiguity.



